\subsubsection{一致有界定理}\label{subsec:uniform-boundedness-theorem}

很容易构造 $[0,1]$ 上的函数序列 $f_n$, 使对每个 $x$, 序列 $f_n(x)$ 有界, 但 $f_n$ 并不一致有界. 例如,
\begin{equation}\label{eq:pointwise-bounded-not-uniform-function-sequence}
f_n(x):=\begin{cases}0,&x\in[0,1/n)\\\dfrac{1}{x},&x\in[1/n,1].\end{cases}
\end{equation}
相比之下, 一致有界定理断言, 对 Banach 空间之间的线性映射, 逐点有界蕴含一致有界.

\begin{Theorem}
\label{thm:uniform-boundedness}
设 $X,Y$ 为 Banach 空间. 设 $T_n\in\mathcal{L}(X,Y)$, $n\in\mathbb{N}$, 满足对每个固定的 $x\in X$, $\|T_nx\|$ 关于 $n$ 一致有界. 那么 $\|T_n\|$ 关于 $n$ 一致有界.
\end{Theorem}

\begin{Corollary}
\label{cor:weakly-bounded-sequence-norm-bound}
设 $Z$ 为 Banach 空间, $x_n$ 为 $Z$ 中的序列, 且对每个 $f\in Z^*$, $f(x_n)$ 关于 $n$ 一致有界. 那么 $\|x_n\|$ 关于 $n$ 一致有界.
\end{Corollary}

在 Theorem~\ref{thm:uniform-boundedness} 中令 $X=Z^*$ 且 $Y=\mathbb{R}$ ($\mathbb{C}$), 就得到这一推论的证明.

\begin{Corollary}
\label{cor:weak-operator-bounds-imply-uniform-bound}
设 $X,Y$ 为 Banach 空间, $T_n\in\mathcal{L}(X,Y)$ 满足对每个 $f\in Y^*$ 及 $x\in X$, $f(T_nx)$ 关于 $n$ 一致有界. 那么 $\|T_n\|$ 关于 $n$ 一致有界.
\end{Corollary}

证明时, 先应用 Corollary~\ref{cor:weakly-bounded-sequence-norm-bound}, 得到 $\|T_nx\|$ 关于 $n$ 一致有界, 再使用 Theorem~\ref{thm:uniform-boundedness}.

Theorem~\ref{thm:uniform-boundedness} 的证明基于 Baire 纲定理. 这里不作这一证明, 而请读者参阅文献. 下面在一个特殊情形中给出简单得多的证明, 即 Corollary~\ref{cor:weakly-bounded-sequence-norm-bound} 中 $Z$ 为 Hilbert 空间的情形. 假设结论不成立, 即 $\|x_n\|$ 不一致有界. 令 $M=\{x_n\mid n\in\mathbb{N}\}$, 并令
\begin{equation}\label{eq:hilbert-unbounded-sequence-functional-supremum}
\alpha(f)=\sup_{n\in\mathbb{N}}|(f,x_n)|
\end{equation}
其中 $f\in Z$. 由于 $M$ 不一致有界, 存在 $y_1\in M$ 及单位向量 $e_1$, 使 $|(e_1,y_1)|\geq1$. 设 $Q$ 为由 $e_1$ 和 $y_1$ 张成的向量空间. 每个 $x_n$ 可分解为 $x_n=x_n^1+x_n^2$, 其中 $x_n^1\in Q$ 且 $x_n^2\in Q^\perp$. 由于 $|(e_1,x_n)|$ 及 $|(y_1,x_n)|$ 一致有界, $x_n^1$ 一致有界, 因而 $x_n^2$ 不一致有界. 所以能找到 $y_2\in M$ 及单位向量 $e_2\in Q^\perp$, 使
\begin{equation}\label{eq:hilbert-unbounded-sequence-second-vector}
|(e_2,y_2)|\geq2(\alpha(e_1)+2).
\end{equation}
归纳地, 可以找到 $y_{n+1}\in M$ 及与 $e_1,e_2,\ldots,e_n;y_1,y_2,\ldots,y_n$ 正交的单位向量 $e_{n+1}$, 使
\begin{equation}\label{eq:hilbert-unbounded-sequence-inductive-vectors}
|(e_{n+1},y_{n+1})|\geq(n+1)\left(\sum_{i=1}^n\dfrac{1}{i}\alpha(e_i)+n+1\right).
\end{equation}
现在令
\begin{equation}\label{eq:hilbert-unbounded-sequence-testing-series}
f=\sum_{i=1}^{\infty}\dfrac{1}{i}e_i.
\end{equation}
该级数收敛, 因为 $e_i$ 为标准正交向量, 且 $\sum_{i=1}^{\infty}i^{-2}<\infty$. 进一步得到
\begin{equation}\label{eq:hilbert-unbounded-sequence-contradiction}
\begin{aligned}
|(f,y_{n+1})|&=\left|\sum_{i=1}^n\dfrac{1}{i}(e_i,y_{n+1})+\dfrac{1}{n+1}(e_{n+1},y_{n+1})\right|\\
&\geq-\sum_{i=1}^n\dfrac{1}{i}\alpha(e_i)+\dfrac{1}{n+1}(n+1)\left(\sum_{i=1}^n\dfrac{1}{i}\alpha(e_i)+n+1\right)\\
&=n+1.
\end{aligned}
\end{equation}
因此序列 $|(f,x_n)|$ 不可能有界, 矛盾.
